\begin{claim}{D.3.4}{Proved}{Large target families and the collapse of Moore wedges}
Let $\kappa$ be an infinite cardinal and $d\geq3$, and put
\[
 Y_\kappa^d=\bigvee_{\alpha<\kappa}P^d(2),\qquad
 \mathcal N_\kappa^d=
       \{X:\Map_*(X,Y_\kappa^d)\simeq *\}.
\]
Then $\mathcal N_\kappa^d$ is closed under desuspension, for
arbitrary pointed source spaces. In fact all its cardinal
parameters collapse:
\[
 \mathcal N_\kappa^d
   =\{X:\Map_*(X,P^d(2))\simeq *\}
   =\{X:\Map_*(X,P^3(2))\simeq *\}.
 \tag{D.3.4.0}
\]
Consequently its cofiber test,
intersected with $R(\loc(\mathbb S/2))$, is a reversible
fundamental localizer $W_\kappa^d$ satisfying
\[
 W_{\mathrm{rev}}\subseteq W_\kappa^d
       \subsetneq R(\loc(\mathbb S/2)),\qquad
 T(W_\kappa^d)=\loc(\mathbb S/2).
 \tag{D.3.4.1}
\]
There is also such a construction using all simply connected
finite-dimensional CW targets with at most $\kappa$ cells.
The infinite Moore wedges give one localizer, independently of
cardinality and dimension. D.3.5 below shows that the broader
target families also give one localizer, namely
$W_S\cap R(\loc(\mathbb S/2))$, independently of cardinality.
Whether either construction differs strictly from
$W_{\mathrm{rev}}$ remains unresolved here.
\end{claim}
\begin{proof}
Fix a source $X$ and again put
$\mathcal R_X=\{Z:\Map_*(X,Z)\simeq *\}$.
We first prove the large-wedge assertion, including its
bottom-dimensional case without a localization convention for
$\RP^2$.

\medskip\noindent
\textbf{A three-factor finite retraction.}
For $A=P^{d-1}(2)$, let $a=d-2\geq1$ and $s=2a+1$.
There is an actual map
\[
 \alpha:S^{2a+1}\longrightarrow A\wedge A,
 \qquad \alpha_*(\iota)=uv+vu\quad\hbox{over }\FF_2.
\]
Here is an elementary verification valid also for $a=1$.
The bottom skeleton of $A\wedge A$ has one cell in degree
$2a$ and two cells in degree $2a+1$. The two attaching degrees
are $2$ and $(-1)^a2$. A unimodular change of the two attaching
spheres changes this row to $(2,0)$. The skeleton is therefore
homotopy equivalent to $P^{2a+1}(2)\vee S^{2a+1}$.
The sphere summand is represented by the integral cycle
$v\otimes u-(-1)^a u\otimes v$; its image modulo two is
$uv+vu$. Include this sphere into the full smash product to
obtain $\alpha$. This is the Moore-space instance of
\cite[Proposition~2.1]{I_ChenWu2013}.

The suspended space $\Sigma A^{\wedge3}$ is simply connected
and has finite $2$-primary integral homology, even when
$A=\RP^2$. The complementary Dynkin--Specht--Wever telescopes
of D.1.8 can therefore be formed on this suspended space:
multiplication by $3$ is a homotopy equivalence and can be
inverted. The same eight-word calculation (D.1.8.4) gives
\[
 \Sigma^{s+1}A=P^{3d-3}(2)
       \quad\hbox{as a homotopy retract of}\quad
       \Sigma A^{\wedge3}.
 \tag{D.3.4.2}
\]
Explicitly, project to the $\beta_3$ telescope and use the
equivalence from $\Sigma^{s+1}A$ induced by
$\Sigma(\id_A\wedge\alpha)$; its two homology images are
$uuv+vuu$ and $uvv+vvu$. The inclusion of the telescope and
the inverse equivalence supply the two retraction maps.
No assertion that $\RP^2$ itself is $2$-local is used.

\medskip\noindent
\textbf{A retraction for the entire infinite wedge.}
Partition the index set into three subsets, each of cardinality
$\kappa$, and write
\[
 Y_\kappa^d\simeq\Sigma U_1\vee\Sigma U_2\vee\Sigma U_3,
 \qquad U_i=\bigvee_{\alpha<\kappa}A.
\]
Use Hilton--Milnor with these three connected spaces as its
three letters. The factor for the basic word $[[1,2],3]$ makes
$\Omega\Sigma(U_1\wedge U_2\wedge U_3)$ a homotopy retract of
$\Omega Y_\kappa^d$. The theorem uses finitely many letters;
there is no finiteness assumption on each letter's CW complex.
Its factors' connectivities tend to infinity with word length,
so the factor inclusions and projections are legitimate.

The smash product here is the wedge of copies of $A^{\wedge3}$
indexed by $\kappa^3$. Choose a diagonal set of $\kappa$ triples,
include those summands, and collapse the others. Wedge together
the actual finite retractions (D.3.4.2) over this diagonal set.
After looping, their composite with the Hilton--Milnor
retraction gives
\[
             \Omega Y_\kappa^{3d-3}
              \quad\hbox{as a retract of}\quad
             \Omega Y_\kappa^d .                   \tag{D.3.4.3}
\]
This construction includes $d=3$: in that case (D.3.4.2) is
the retract $P^6(2)$ of $\Sigma(\RP^2)^{\wedge3}$.
The use of three whole wedges, followed by an actual wedge
projection, avoids replacing an infinite collection of
same-connectivity Hilton--Milnor factors by a product.

Suppose $\Omega Y_\kappa^d\in\mathcal R_X$, and set
$D=3d-3>d$ and $K=Y_\kappa^{D-1}$.
Equation (D.3.4.3) gives $\Omega\Sigma K\in\mathcal R_X$.
Since $\kappa+\kappa=\kappa$, the space $\Sigma K$ is
homotopy equivalent to $\Sigma K\vee\Sigma K$.
The two-letter Hilton--Milnor theorem now puts
\[
                     \Omega\Sigma K^{\wedge q}\in\mathcal R_X
                     \qquad(q\geq1).
\]
Apply the two-copy argument proved in D.1.8, treating the
whole simply connected space $K$ as one letter. For every
finite wedge $V$ of copies of $K$, its iterated James
splittings have finite multiplicities at each smash weight.
The tails become arbitrarily highly connected, regardless of
the cardinality of each weight. Thus all the positive iterates
$(\Omega\Sigma)^{r+1}V$ belong to $\mathcal R_X$, and the
convergent simply connected resolution gives
$V\in\mathcal R_X$.
Barratt's finite cone cubes lower this whole finite-wedge
family from $Y_\kappa^{D-1}$ to $Y_\kappa^d$.
We have proved
\[
 \Map_*(X,\Omega Y_\kappa^d)\simeq *
       \quad\Longrightarrow\quad
 \Map_*(X,Y_\kappa^d)\simeq *.
\]
All inverse limits are in the targets. This proves the
arbitrary-source desuspension assertion.

\medskip\noindent
\textbf{Why the Moore cardinal parameter collapses.}
We strengthen the result just proved. Suppose first that $X$
is countable and $P^d(2)\in\mathcal R_X$. The proof of D.1.8,
including its iteration and finite cone cubes, puts every
finite bouquet of copies of $P^d(2)$ in $\mathcal R_X$.
Write $A=P^{d-1}(2)$ and $F_q=\Sigma A^{\wedge q}$ for
$q\geq1$. The finite-letter Hilton--Milnor theorem therefore
gives $\Omega F_q\in\mathcal R_X$ for every $q$.
Each $F_q$ is a finite simply connected complex with finite
$2$-primary integral homology. Its homotopy groups are finite
$2$-groups, by the Serre-class version of the finiteness
theorem. Also its connectivity tends to infinity with $q$.

Choose an exhaustion $X=\bigcup_i X_i$ by finite pointed
subcomplexes. It is not necessary that $X$ have finite skeleta.
For fixed $q$, the groups
\[
                 [X_i,\Omega F_q]_*
\]
are finite. The Milnor sequence and the vanishing of $\lim^1$
for towers of finite groups identify their inverse limit with
$[X,\Omega F_q]_*=0$. A tower of finite groups with zero
inverse limit is pro-zero: at any fixed stage its descending
images stabilize, and the stabilized image is the image of
the inverse limit, so is trivial. The same assertion holds
with $X_i$ replaced by any fixed suspension $\Sigma^rX_i$.

The Hilton--Milnor decomposition for a countable bouquet,
obtained compatibly from its finite subbouquets, is the weak
product of copies of $\Omega F_q$, indexed by basic words of
length $q$. For any finite $X_i$, maps to this weak product
factor through a finite subproduct. Thus it identifies the
restriction tower $[X_i,\Omega Y_{\aleph_0}^d]_*$, as a tower
of pointed sets, with the weak products of the factor towers.
For fixed $i$, only finitely many values of $q$ can be visible,
by the dimension of $X_i$ and the connectivity of $F_q$.
Choose $j\geq i$ that kills the image of every one of those
finitely many factor types. It then kills every copy of each
type, regardless of its word label. Consequently
\[
 [X_j,\Omega Y_{\aleph_0}^d]_*\longrightarrow
 [X_i,\Omega Y_{\aleph_0}^d]_*
             \quad\text{is the zero map}.
 \tag{D.3.4.4}
\]
This step uses the weak product itself, not an identification
of it with the full product. The same proof applies to all
the suspensions $\Sigma^rX_i$. These pro-zero towers have
zero inverse limit and zero $\lim^1$, so their Milnor
sequences give
\[
             \Map_*(X,\Omega Y_{\aleph_0}^d)\simeq *.
\]
The desuspension result already proved makes
$Y_{\aleph_0}^d\in\mathcal R_X$. Every map from a suspension
of countable $X$ to $Y_\kappa^d$ factors through a countable
subbouquet: the image of each closed source cell lies in a
finite target subcomplex. That subbouquet is a retract of
$Y_{\aleph_0}^d$. Hence $Y_\kappa^d\in\mathcal R_X$ for all
$\kappa$. This is the uniform finite-factor-type argument
behind \cite[Proposition~6]{I_Strom2012}, with a finite-cell
exhaustion and finite homotopy groups replacing its
finite-type and rational hypotheses.

To pass from countable to arbitrary $X$, apply the bounded
extension lemma of G.7 to a finite simplicial Moore model
$K$ and its Bousfield--Kan $\FF_2$-resolution. Exactly the
compact model construction in G.7 makes
\[
                  Z=\operatorname{Tot}(R^{\bullet+1}K)
\]
a Kan complex with compact metrizable levels. It represents
the $2$-completion of $P^d(2)$. Since the latter space is
simply connected and all its homotopy groups are finite
$2$-groups, the completion map is a weak equivalence. Apply
the lemma simultaneously to $Z,\Omega Z,\Omega^2 Z,\ldots$.
It says that every countable simplicial subcomplex of a
$P^d(2)$-null $X$ is contained in a countable $P^d(2)$-null
subcomplex. These countable null subcomplexes form a filtered
system with homotopy colimit $X$. Each is $Y_\kappa^d$-null
by the countable argument. Nullity is closed under pointed
homotopy colimits, so $X$ is $Y_\kappa^d$-null as well.
The reverse implication uses the retraction onto one target
summand. Equality of the dimensions follows from D.1.8.
This proves (D.3.4.0) for arbitrary sources.

\medskip\noindent
\textbf{The family with a cell-cardinality bound.}
Let $\mathcal F_\kappa$ consist of representatives of the simply
connected finite-dimensional CW complexes with at most
$\kappa$ cells, and put
\[
 \mathcal N(\mathcal F_\kappa)=
       \{X:\Map_*(X,Z)\simeq *\text{ for all }Z\in\mathcal F_\kappa\}.
\]
Suspension, finite wedges and finite smash products preserve
this target family up to weak equivalence, since $\kappa$ is
infinite. Suppose $\Sigma X\in\mathcal N(\mathcal F_\kappa)$.
Then $\Omega Z\in\mathcal R_X$ for every
$Z\in\mathcal F_\kappa$. For fixed $Z$ in the family, each
finite weight truncation of
\[
 \Sigma(\Omega\Sigma)^rZ
     \simeq\bigvee_{m\geq1}\bigvee^{c_{r,m}}\Sigma Z^{\wedge m}
\]
is finite-dimensional and has at most $\kappa$ cells. Its loop
space is therefore in $\mathcal R_X$. Increasing connectivity
of the tails puts $(\Omega\Sigma)^{r+1}Z$ in that class; the
convergent resolution puts $Z$ there. This holds for every
$Z\in\mathcal F_\kappa$, proving desuspension closure of
$\mathcal N(\mathcal F_\kappa)$.

\medskip\noindent
\textbf{The stable part and the precise remaining question.}
Both null classes are closed under pointed homotopy colimits.
Desuspension and the mapping-space fibre sequence give all
three cofiber implications, as in A.3. Their cofiber tests
therefore define reversible fundamental localizers. Intersect
each with $R(\mathcal L)$, where $\mathcal L=\loc(\mathbb S/2)$.
Miller's theorem \cite[Theorem~A]{Miller1984} puts
$B\ZZ/2$ in each null class: all the targets used are
finite-dimensional, including the possibly enormous Moore
wedge. Its reduced suspension spectrum belongs to $\mathcal L$.
Thus the intersection contains $W_{\mathrm{rev}}$, and its
stable part contains $T(W_{\mathrm{rev}})=\mathcal L$.
The opposite containment follows from the intersection with
$R(\mathcal L)$, proving equality of stable parts.

The finite space $P^d(2)$ is in the upper endpoint's aspherical
class and is outside both new null classes. For the single
large wedge use its inclusion as one wedge summand, followed
by the retraction onto that summand; for
$\mathcal F_\kappa$ use $P^d(2)$ itself as target. These maps
are not null, proving the strict upper inclusion in (D.3.4.1).

For the family $\mathcal F_\kappa$, larger cardinal bounds
initially yield smaller or equal null classes. The complete
countable-source and compactness argument in D.3.5 shows that
all these interval localizers are equal to
$W_S\cap R(\loc(\mathbb S/2))$. Thus cardinality cannot
supply a separator within this family. Its distinctness
from $W_{\mathrm{rev}}$ remains the relevant unresolved issue.
\end{proof}
